\E{acceleratrice}{əkˌsɛləˈreɪtrɪs}{n.}{Fr. \la{accélératrice}, fem. agentive f. L. \la{accelerare} to hasten} \sn{1} A machine, institution or person whose function is to quicken other processes, as a particle-accelerator, an incubator of ventures, or a well-placed patron. \sn{2} \lbl{fig.} A war, plague or rivalry considered as an involuntary hastener of invention. \ex{Necessity is not the mother of invention but its acceleratrice.} \src{Hom. Grad.} iii. Cf. \xr{amplificatrix}.

\E{actinophage}{ækˈtɪnəfeɪdʒ}{n.}{Gk. ἀκτίς, ἀκτῖνος ray + -φάγος eating} A ray-eater; a solar collector, or any surface laid out to drink the light of the sun. \ex{The deserts were tiled with actinophages, and the night-cities glowed on daylight banked.} \src{Tract. Cop.} viii. --- \dv{actinophagy} \textit{n.}; \dv{actinophagous} \textit{adj.} Cf. \xr{heliophage}.

\E{adynamocrat}{əˈdaɪnəməʊkræt}{n.}{Gk. ἀδύναμος powerless + κράτος rule} An official whose office consists in rendering others powerless; the administrator of permits not granted and of connections to the grid indefinitely deferred. --- \dv{adynamocracy} \textit{n.} Cf. \xr{bradycrat}.

\E{aethergy}{ˈiːθədʒi}{n.}{Gk. αἰθήρ the upper air + ἔργον work} Energy carried through the air without wires; by extension, the whole art of radio, beam and broadcast, regarded as a sending of power through the heavens. --- \dv{aethergic} \textit{adj.}

\E{agonarch}{ˈæɡənɑːk}{n.}{Gk. ἀγών contest + ἀρχός leader} The leader in a race of technologies; the laboratory or nation that stands first upon the \xr{kairodrome}, and is consequently first to be asked to pause. \ex{Every agonarch is a moratorian's first target.} \src{Sat. Decel.} x. --- \dv{agonarchy} \textit{n.}

\E{akinetist}{əˈkɪnɪtɪst}{n.}{Gk. ἀκίνητος unmoved, not to be moved} One who holds that the safest policy is no motion at all, and who regards each proposal as presumptively guilty. Distinguished from the \xr{inertist} by being articulate about it. --- \dv{akinetism} \textit{n.}

\E{alethotachy}{əˈliːθəʊˌtæki}{n.}{Gk. ἀλήθεια truth + τάχος speed} The speed at which a truth propagates through the Agora; proverbially slower than that of its refutation, and slower still than that of the meme mocking both. \ex{A lie is round the Agora twice before alethotachy has laced its boots.} \src{Agora}

\E{amechanist}{əˈmɛkənɪst}{n.}{Gk. ἀμήχανος without means, helpless} One who refuses machines on principle and is thereby rendered helpless; the term preserves the Greek sense, in which to be without devices is to be without resource. --- \dv{amechany} \textit{n.} Cf. \xr{mechanoclast}.

\E{anergist}{ˈænədʒɪst}{n.}{Gk. ἀ- \lbl{priv.} + ἔργον work} One who would do without energy, or with markedly less of it; who counts a darkened city a moral achievement. Opp. \xr{energolatry}. --- \dv{anergism} \textit{n.}; \dv{anergic} \textit{adj.}

\E{anthropergy}{ænˈθrɒpədʒi}{n.}{Gk. ἄνθρωπος man + ἔργον work} The energy of human muscle, considered as a historical fuel; the measure by which pre-industrial civilizations were reckoned, and beyond which the engine first lifted mankind. \ex{Every pyramid was built of anthropergy, and every pyramid was small.} \src{Lib. Ascens.} iv.

\E{antitachist}{ˌæntiˈtækɪst}{n.}{Gk. ἀντί against + τάχος speed} An avowed opponent of speed in technological affairs; less euphemistic than the \xr{decelerist}, and so more respected by the accelerationists, who prefer an honest enemy. --- \dv{antitachism} \textit{n.}

\E{apragmocrat}{əˈpræɡməʊkræt}{n.}{Gk. ἀπράγμων inactive, keeping out of affairs + κράτος rule} A ruler whose virtue is held to consist in doing nothing; an official praised for every disaster that did not occur, and never blamed for any benefit that did not. --- \dv{apragmocracy} \textit{n.} Cf. \xr{acediocracy}.

\E{argyrogenesis}{ˌɑːdʒɪrəʊˈdʒɛnɪsɪs}{n.}{Gk. ἄργυρος silver, money + γένεσις birth} The creation of money; \sn{1} the issuing of currency; \sn{2} \lbl{laudatory} the generation of capital by new enterprise, as distinguished from its extraction. Cf. \xr{plutogenesis}.

\E{astrergy}{ˈæstrədʒi}{n.}{Gk. ἄστρον star + ἔργον work} The energy of the stars, considered as the inheritance awaiting any civilization bold enough to collect it. \ex{The galaxy wastes its astrergy into the dark, and calls the waste beauty.} \src{Lib. Ascens.} vi. --- \dv{astrergic} \textit{adj.}

\E{asymptotarian}{əˌsɪmptəˈtɛəriən}{n.}{Gk. ἀσύμπτωτος not falling together + -arian} One who maintains that every curve of progress must flatten, and who greets each record as proof the plateau is near. Opp. \xr{hyperbolist}. --- \dv{asymptotarianism} \textit{n.} Cf. \xr{asymptophobia}.

\E{autarkergy}{ɔːˈtɑːkədʒi}{n.}{Gk. αὐτάρκεια self-sufficiency + ἔργον work} Self-sufficiency in energy; the condition of a city, nation or colony that generates all the power it consumes, and so answers to no supplier. --- \dv{autarkergic} \textit{adj.}

\E{axiotachy}{ˈæksiəʊˌtæki}{n.}{Gk. ἄξιος worthy + τάχος speed} The value of speed; the doctrine that a benefit delivered sooner is a greater benefit, and that delay must therefore be priced in every calculation of good. \ex{Axiotachy is the discount rate of the righteous.} \src{Cat. Acc.} 9.

\E{baryergy}{ˈbærɪɜːdʒi}{n.}{Gk. βαρύς heavy + ἔργον work} Energy drawn from weight and fall, as from dams, tides and gravitational storage; the patient power of things coming down. --- \dv{baryergic} \textit{adj.}

\E{bathmoclast}{ˈbæθməʊklæst}{n.}{Gk. βαθμός step + κλάστης breaker} One who would kick away the ladder; who, having climbed by cheap energy and free enquiry, would deny both to those beneath. \ex{The bathmoclast preaches sufficiency from the top stair.} \src{Sat. Decel.} vii. Cf. \xr{anabathmism}.

\E{bathyergy}{ˈbæθɪɜːdʒi}{n.}{Gk. βαθύς deep + ἔργον work} Deep energy; the heat of the planet's interior, tapped by drilling, and held by the anabathmist to be the most neglected fortune on earth. --- \dv{bathyergic} \textit{adj.}

\E{bibliostasis}{ˌbɪbliəʊˈsteɪsɪs}{n.}{Gk. βιβλίον paper, document + στάσις standing still} The arrest of a project by the weight of its documents; a stagnation produced not by refusal but by the accumulation of required reports. \ex{The reactor died of bibliostasis at ten thousand pages.} \src{Tract. Cop.} iv. Cf. \xr{grammatocracy}.

\E{biodromy}{baɪˈɒdrəmi}{n.}{Gk. βίος life + δρόμος course} The course of a life considered as a race-track that may be lengthened; the span of years, regarded by the longevist as a parameter rather than a fate. --- \dv{biodromic} \textit{adj.} Cf. \xr{macrobiotarch}.

\E{blastotachy}{ˈblæstəʊˌtæki}{n.}{Gk. βλαστός sprout, shoot + τάχος speed} The rate at which new ventures sprout in a society; the vernal vigour of a startup ecology. \ex{Cut the permits and watch the blastotachy.} \src{Ep. Fabr.} iv. 2.

\E{boulomancy}{ˈbuːləʊˌmænsi}{n.}{Gk. βουλή council + μαντεία divination} The art of predicting what a committee will decide; held to be easier than any other branch of prophecy, since the answer is almost always ‘later’. --- \dv{boulomancer} \textit{n.} Cf. \xr{synedriasm}.

\E{brachylogism}{brəˈkɪlədʒɪzəm}{n.}{Gk. βραχυλογία brevity of speech} An argument compressed to the length of a single post; the characteristic unit of reasoning on the Agora, praised by the Laconian and deplored by the essayist. --- \dv{brachylogist} \textit{n.} Cf. \xr{memetaxis}.

\E{bradyclast}{ˈbrædɪklæst}{n.}{Gk. βραδύς slow + κλάστης breaker} A breaker of slowness; a reformer who dismantles queues, waiting-periods and review stages. Opp. \xr{bradycrat}. --- \dv{bradyclasm} \textit{n.}

\E{bradygenic}{ˌbrædɪˈdʒɛnɪk}{adj.}{Gk. βραδύς slow + -γενής born of} Generating delay; said of rules, forms and meetings whose principal product is time lost. \ex{a bradygenic statute of four hundred clauses}.

\E{bradynomy}{brəˈdɪnəmi}{n.}{Gk. βραδύς slow + νόμος law} The body of law that governs the pace of things rather than their content: waiting-periods, comment-windows, moratoria. --- \dv{bradynomic} \textit{adj.} Cf. \xr{nomophragm}.

\E{bradyplasty}{ˈbrædɪplæsti}{n.}{Gk. βραδύς slow + πλάσσειν to mould} The reshaping of a bold proposal until it moves slowly enough to pass; the art of the amendment. \ex{The bill survived bradyplasty and emerged unrecognisable.} \src{Min. Pausiar.}

\E{bradypolitan}{ˌbrædɪˈpɒlɪtən}{n. \& adj.}{Gk. βραδύς slow + πολίτης citizen} A citizen of a slow city; one at ease in a place where nothing has been built in a generation. Opp. \xr{nautopolis}.

\E{brontergy}{ˈbrɒntədʒi}{n.}{Gk. βροντή thunder + ἔργον work} The energy of the storm; electricity in its raw and celestial state, before it was tamed into wires. \ex{Franklin stole brontergy with a kite, and no committee was consulted.} \src{Chron. Fabr.}

\E{brotolatry}{brəʊˈtɒlətri}{n.}{Gk. βροτός mortal + λατρεία worship} The worship of mortality; the doctrine that death gives life its meaning, and that its abolition would be an impiety. Opp. \xr{macrobiotarch}. --- \dv{brotolater} \textit{n.}

\E{bucolomania}{ˌbjuːkələʊˈmeɪniə}{n.}{Gk. βουκολικός pastoral + μανία madness} A mania for the rural idyll; the longing, generally felt in cities and typed upon machines, for a pre-industrial life its sufferers have never lived. --- \dv{bucolomaniac} \textit{n.} Cf. \xr{retrotropism}.

\E{cacergy}{ˈkækədʒi}{n.}{Gk. κακός bad + ἔργον work} Wasted energy; heat dissipated without work done, as in a meeting. \ex{Half the economy runs on cacergy.} \src{Brev. Therm.} ix. Opp. \xr{eleutherergy}.

\E{cathedrocrat}{kəˈθiːdrəkræt}{n.}{Gk. καθέδρα seat, chair + κράτος rule} One who governs from the armchair; a theorist of technologies never built, whose authority is the more secure for never having been tested. --- \dv{cathedrocracy} \textit{n.} Cf. \xr{logomachist}.

\E{celerarch}{ˈsɛlərɑːk}{n.}{L. \la{celer} swift + Gk. ἀρχός ruler} A master of speed; a manager whose distinction is that things in their charge are finished sooner than anyone believed possible. --- \dv{celerarchy} \textit{n.} Cf. \xr{tachycrat}.

\E{centuplator}{ˈsɛntjʊpleɪtə}{n.}{L. \la{centuplex} hundredfold + agentive -ator} One who seeks hundredfold returns; an investor in the improbable, who expects most ventures to die and one to pay for all. --- \dv{centuplation} \textit{n.}

\E{chemergy}{ˈkɛmədʒi}{n.}{Gk. χημεία chemistry + ἔργον work} Chemical energy; the power of the flame and the battery, regarded by the anabathmist as a provisional stage before the nucleus. Cf. \xr{ultrapyric}.

\E{chiliodrome}{ˈkɪliədrəʊm}{n.}{Gk. χίλιοι a thousand + δρόμος course} The thousand-year race; the course of a civilization reckoned in millennia, along which present quarrels appear as stumbles at the starting line. \ex{Run the chiliodrome as if it were a sprint.} \src{Hom. Grad.} xi.

\E{chronoclast}{ˈkrɒnəʊklæst}{n.}{Gk. χρόνος time + κλάστης breaker} A breaker of schedules; one who finishes years ahead of the plan, to the embarrassment of the planners. --- \dv{chronoclastic} \textit{adj.}

\E{chronotomy}{krəˈnɒtəmi}{n.}{Gk. χρόνος time + τομή a cutting} The cutting of time; the compression of a projected timeline, as when a century's research is expected in a decade. --- \dv{chronotomist} \textit{n.} Cf. \xr{brachychronist}.

\E{cogitarium}{ˌkɒdʒɪˈtɛəriəm}{n., pl. -ia}{L. \la{cogitare} to think + -arium place} A place of thinking; a cluster of machines assembled for the training of minds. \ex{They built the cogitarium beside the river, for the cooling.} \src{Chron. Fabr.} Cf. \xr{logisterion}.

\E{colonifer}{kəˈlɒnɪfə}{n.}{L. \la{colonia} settlement + \la{-fer} bearing} A colony-bearer; a ship or launch-vehicle that carries settlers off-world. --- \dv{coloniferous} \textit{adj.} Cf. \xr{apoikiast}.

\E{compendiarch}{kəmˈpɛndiɑːk}{n.}{L. \la{compendium} a short cut + Gk. ἀρχός ruler} A ruler by shortcut; one who finds the unbureaucratic path and dares to take it. --- \dv{compendiarchy} \textit{n.}

\E{conflagratrice}{ˌkɒnfləˈɡreɪtrɪs}{n.}{Fr. \la{conflagratrice}, f. L. \la{conflagrare} to burn up} That which sets things ablaze; an idea, invention or post that ignites a whole field at once. \ex{The paper was a conflagratrice; within the year every lab was burning.} \src{Agora}

\E{consiliophage}{kənˈsɪliəfeɪdʒ}{n.}{L. \la{consilium} council + Gk. -φάγος eating} An eater of councils; a reformer who abolishes committees, or a technology that renders them obsolete. Cf. \xr{bouleutophobia}.

\E{copiogenesis}{ˌkəʊpiəʊˈdʒɛnɪsɪs}{n.}{L. \la{copia} abundance + Gk. γένεσις birth} The creation of abundance; the process by which a scarce thing is made cheap. \ex{Copiogenesis ends more wars than treaties do.} \src{Tract. Cop.} i. Cf. \xr{copiarch}.

\E{cosmagogue}{ˈkɒzməɡɒɡ}{n.}{Gk. κόσμος world + ἀγωγός leading} A leader into the cosmos; one who conducts a people off-world. --- \dv{cosmagogy} \textit{n.}

\E{crescentism}{krɪˈsɛntɪzəm}{n.}{L. \la{crescens} growing; Fr. \la{crescentisme}} The doctrine of perpetual growth; the moral preference for the waxing over the full moon. --- \dv{crescentist} \textit{n.} Cf. \xr{auxetist}.

\E{cryptostasis}{ˌkrɪptəʊˈsteɪsɪs}{n.}{Gk. κρυπτός hidden + στάσις standing still} Hidden stagnation; the condition of a society that appears busy but builds nothing new. Cf. \xr{kenotachy}.
